\chapter{Calculus}

\section{Exponents and logarithms}
The exponential function is
\[
e^x = 1 + x + \frac{x^2}{2} + \frac{x^3}{3!} + \cdots
\]
The function $\log$ is the inverse of the exponential:
\[
\log e^x = x
\]


For small $x$, truncating after the linear term gives the useful approximation
\[
e^x \approx 1+x \qquad (x \text{ small}).
\]
Since $\log$ is the inverse of $e^x$, the same approximation, read the other way, gives
\[
\log(1+x) \approx x \qquad (x \text{ small}).
\]

\subsection{Identities}
\begin{align*}
  \left (e^x\right )^y &= e^{xy} \\
  \log x^y &= y \log x \\
  e^{x+y} &= e^x e^y.
\end{align*}
%% Derivative of the exponential:
%% \[
%% \frac{d}{dx}e^x = e^x, \qquad \frac{d}{dx}e^{kx} = k\,e^{kx}.
%% \]

\section{The derivative as a linear approximation}
\[
y+\Delta y = y(x+\Delta x) \approx y(x) + y'(x)\Delta x
\quad\Longrightarrow\quad
\Delta y = y'(x)\Delta x.
\]
$y'(x)$ is called \emph{the derivative}; it is also written $\dfrac{dy}{dx}$.

\subsection{Multiplication by a scalar}
What is the derivative of $z(x) = y(\alpha x)$? By the linear approximation:
\[
z + \Delta z \approx z(x) + z'(x) \Delta x
\]
But, on the other hand, by the definition of $z$:
\begin{align*}
  z + \Delta z &= z(x+\Delta x)\\
  &= y(\alpha x + \alpha \Delta x) \\
  &=y(\alpha x) + y'(\alpha x) \alpha \Delta x \\
  &=z(x) + y'(\alpha x) \cdot  \alpha \Delta x
\end{align*}
Therefore
\[
z'(x) = \alpha y'(\alpha x)  
\]
We will see later what this means for change of units. In the
meantime, note a specific example: $(e^{\alpha x})' = \alpha e^{\alpha x}$.
\subsection{Other identities}
\begin{align*}
  &\text{Sum rule:}&& (y+z)'(x) = y'(x) + z'(x) \\
  &\text{Product / Leibniz rule:}&&   (yz)'(x) = y'(x)z(x) + y(x)z'(x) \\
  &\text{Chain rule:}&&   \left (y(z(x))\right )' = y'(z(x)) \cdot z'(x)
\end{align*}



%% Our main unit of measurement is time, denoted $t$, measured in (say)
%% days. Let $I_d(t)$ be the number of infected individuals in a population
%% after $t$ days have passed. Then $\dot{I}_d$ is the change in the number
%% of infections per day, i.e.~the rate of new infections. How many new
%% infections can we expect in a week? Clearly that number is $7
%% \dot{I}_d$! But why? What has happened here?

%% Denote $\tau = t / 7$ --- the number of \emph{weeks} passed. So
%% the number of infected individuals after $\tau$ weeks is
%% $I_w(\tau) = I_d(7\tau)$. Then the rate of new infections per
%% \emph{week} is:
%% \begin{align*}
%%   \frac{d I_{w}}{d\tau} &=  7 \frac{d I_d}{d t} 
%%  \end{align*}

%% So we conclude: converting time to units $7$ times coarser multiplies
%% the rate by $7$. In general, measuring $t$ in units $k$ times coarser
%% (so $\tau = t/k$) multiplies any rate by $k$:
%% \[
%% \frac{dI_\star}{d\tau} = k\,\frac{dI}{dt}.
%% \]


\section{The Fundamental Theorem of Calculus}
\[
\int_a^b y'(x)\,dx = y(b) - y(a).
\]

\subsection{Why this works}
Partition $[a,b]$ into $n$ steps $a=x_0<x_1<\cdots<x_n=b$, each of
width $\Delta x$. By the linear approximation of Section~2, each step
changes $y$ by $\Delta y_i = y(x_{i+1})-y(x_i)\approx y'(x_i)\Delta x$. Summing the
steps telescopes:
\begin{align*}
  \sum_{i=0}^{n-1} y'(x_i)\,\Delta x & \approx \sum_{i=0}^{n-1} \Delta y_i \\
  &= \sum_{i=0}^{n-1}\underbrace{\big(y(x_{i+1})-y(x_i)\big)}_{\Delta y_i} \\
  &= \big(y(x_1)-y(x_0)\big) + \big(y(x_2)-y(x_1)\big) + \cdots +
  \big(y(x_n)-y(x_{n-1})\big)\\
  &= y(x_n)-y(x_0) \\
  &= y(b)-y(a).
\end{align*}

\begin{example}
Since $\frac{d}{ds}\left(-\frac{1}{\lambda}e^{-\lambda s}\right) = e^{-\lambda s}$, for $\lambda>0$ (so that $e^{-\lambda s}\to 0$ as $s\to\infty$),
\[
\int_t^\infty e^{-\lambda s}\,ds
= \left[-\frac{1}{\lambda}e^{-\lambda s}\right]_t^\infty
= \left(-\frac{1}{\lambda}e^{-\lambda \cdot \infty} + \frac{1}{\lambda}e^{-\lambda t}\right)
= \frac{1}{\lambda}e^{-\lambda t}.
\]
\end{example}
